\newpage
\appendix
\renewcommand{\thesection}{Annexe \Alph{section}.}
\renewcommand{\thesubsection}{\Alph{section} \Roman{subsection}.}
\titleformat{\subsection}
  {\normalfont\normalsize\bfseries}{\thesubsection}{1em}{}
\titleformat{\section}
  {\normalfont\Large\bfseries\filcenter}{\thesection}{1em}{}

\begin{multicols}{2}

\section{Calculs} \label{sec:calculs}

\begin{equation}
    R_{ik} = 8\pi\left( T_{ik} - \tfrac{1}{2}g_{ik}T \right) - \Lambda g_{ik}
\end{equation}

\lipsum[10-12]

\end{multicols}

\section{Tableaux} \label{sec:tableaux}

\begin{table*}[ht]
\centering
\caption{\textbf{Données brutes complètes}}
\label{tab:donnees}
\begin{tabular}{cc|ccc}
\toprule
\toprule
Série & Essai & $U$ (\si{\volt}) & $B$ (\si{\milli\tesla}) & $d$ (\si{\milli\meter}) \\
\midrule
1 & 1 & 500 & 12.3 & 4.7 \\
1 & 2 & 500 & 12.3 & 4.8 \\
2 & 1 & 500 & 15.0 & 5.9 \\
2 & 2 & 500 & 15.0 & 6.0 \\
\midrule
\midrule
\end{tabular}
\end{table*}

\section{Code Python} \label{sec:python}
\begin{python}
import numpy as np

def charge_specifique(U, B, d):
    """Calcule le rapport q/m."""
    return 2 * U / (B**2 * d**2)

q_m = charge_specifique(500, 12.3e-3, 4.75e-3)
print(f"q/m = {q_m:.3e} C/kg")
\end{python}
